\documentclass{article}

% Recommended, but optional, packages for figures and better typesetting:
\usepackage{microtype}
\usepackage{graphicx}
\usepackage{subfigure}
\usepackage{booktabs} % for professional tables

% hyperref makes hyperlinks in the resulting PDF.
\usepackage{hyperref}

% Attempt to make hyperref and algorithmic work together better:
\newcommand{\theHalgorithm}{\arabic{algorithm}}

% Use the following line for the initial blind version submitted for review:
\usepackage[accepted]{icml2025}

% If accepted, instead use the following line for the camera-ready submission:
% \usepackage[accepted]{icml2025}

% For theorems, math, and such
\usepackage{amsmath}
\usepackage{amssymb}
\usepackage{mathtools}
\usepackage{amsthm}

% Added packages from your original document
\usepackage{multirow}
\usepackage{xcolor}
\usepackage{colortbl}
\usepackage{array}

% ---- colour helpers ----
\definecolor{bestgreen}{RGB}{210,240,210}
\definecolor{warnred}{RGB}{255,210,210}
\definecolor{neutralgray}{RGB}{238,238,238}

% The \icmltitle you define below is probably too long as a header.
% Therefore, a short form for the running title is supplied here:
\icmltitlerunning{SAE-Guided Modality-Specific Unlearning in Vision-Language Models}

\begin{document}

\twocolumn[
\icmltitle{SAE-Guided Modality-Specific Unlearning\\[4pt]
           in Vision-Language Models}

% List of affiliations: The first argument should be a (short)
% identifier you will use later to specify author affiliations
\begin{icmlauthorlist}
\icmlauthor{Aditya Bhatia}{inst1}
\icmlauthor{Aryan Chaudhary}{inst1}
\icmlauthor{Ayush Kumar Gupta}{inst1}
\icmlauthor{Mehul Chowksi}{inst1}
% If you had a third author at a different place:
% \icmlauthor{Jane Doe}{inst2}
\end{icmlauthorlist}

% Define what those short keys (inst1, inst2) actually stand for:
\icmlaffiliation{inst1}{IIIT Hyderabad}
% \icmlaffiliation{inst2}{Company Name, City, Country}

% Add corresponding author emails:
\icmlcorrespondingauthor{Aditya Bhatia}{aditya.bhatia@reserach.iiit.ac.in}
\icmlcorrespondingauthor{Aryan Chaudhary}{aryan.chaudhary@reserach.iiit.ac.in}
\icmlcorrespondingauthor{Ayush Kumar Gupta}{ayush.gupta@reserach.iiit.ac.in}
\icmlcorrespondingauthor{Mehul Chowksi}{mehul.chowksi@reserach.iiit.ac.in}


\icmlcorrespondingauthor{Anonymous Authors}{anonymous@example.com}

\icmlkeywords{Machine Learning, ICML, Vision-Language Models, Unlearning, Sparse Autoencoders}

\vskip 0.3in
]

% this must go after the closing bracket ] following \twocolumn[ ...
\printAffiliationsAndNotice{}  % leave blank if no need to mention equal contribution

\begin{abstract}
We investigate whether Sparse Autoencoders (SAEs) can enable
\emph{surgical, weight-free} concept unlearning inside a
Vision-Language Model (VLM). Concretely, we target the
\textit{zebra} concept in LLaVA-NeXT-8B and compare three
intervention strategies---zero ablation, negative clamping, and
additive decoder-weight steering---applied independently to the
language model (LM) and the CLIP vision encoder.
On the language side, additive SAE steering at Layer~24 with the
top-20 zebra features and suppression coefficient $\alpha\!=\!-4.4$
achieves an \textbf{85.7-point} reduction in Forget Accuracy
(FA: $98.7\%\!\to\!13.0\%$), outperforming Language Gradient Ascent
($35.4$-point reduction) while maintaining Retain Accuracy at $50.0\%$.
On the vision side, our custom L1-ReLU SAE trained on CLIP Layer~18
achieves complete zebra suppression ($\mathrm{FA}=0\%$) with as few as
20 features at $\alpha\!=\!-2.0$.
We characterise a \emph{Goldilocks zone} for the suppression
coefficient, document cross-modal knowledge leakage, and provide
geometric evidence of feature superposition as the mechanism
underlying collateral damage to related concepts (e.g.\ crosswalk
recognition).
\end{abstract}

% =====================================================================
\section{Introduction}
\label{sec:intro}

Machine unlearning—the removal of specific knowledge from a trained
model without full retraining—has become a pressing concern for large
generative models, where factual errors, copyrighted content, or harmful
concepts may need to be excised post-deployment. Gradient-based
approaches such as Gradient Ascent (GA) and Negative Preference
Optimisation (NPO) modify model weights globally, often degrading
unrelated capabilities and leaving no interpretable audit trail of what
was changed.

Sparse Autoencoders (SAEs)~\citep{elhage2022superposition,
anthropic2023monosemanticity} offer an alternative: they decompose dense
residual-stream activations into sparse, near-monosemantic features, each
corresponding approximately to one human-interpretable concept. This
structure makes it possible to \emph{surgically} suppress a target concept
by steering the relevant decoder directions---without touching model
weights, and with full reversibility.

SAEs were initially developed for language models. More recently, they
have been extended to vision encoders: the ExplainableML/sae-for-vlm
work \citep{pach2025saevlm} trained SAEs on CLIP ViT-L/14-336 and
showed that intervening on visual features directly changes the text
output of LLaVA. Matryoshka SAEs \citep{wolodja2025msae} and multimodal
SAEs \citep{evolvingLMMs2025} expanded coverage further.

\paragraph{Our contribution.}
We ask: \emph{can SAE-based steering achieve meaningful concept
unlearning in a VLM, and what happens when only one modality is
steered?} Specifically, we:
\begin{itemize}\setlength\itemsep{2pt}
  \item Compare three intervention strategies---zero ablation, negative
        clamping, and additive steering---for LM-side unlearning,
        showing additive steering decisively outperforms the others and
        Gradient Ascent.
  \item Train custom L1-ReLU SAEs on CLIP ViT-L/14-336, identify
        Layer~18 as the optimal intervention depth after a systematic
        layer sweep, and evaluate vision-side unlearning.
  \item Characterise the \emph{Goldilocks zone} for the suppression
        coefficient $\alpha$, beyond which token collapse occurs.
  \item Provide geometric evidence of \emph{feature superposition} to
        explain why crosswalk recognition is partially damaged by zebra
        suppression.
  \item Measure cross-modal knowledge leakage: after language-only SAE
        steering, Language Leakage drops from $96.0\%$ to $60.0\%$
        while Retain Accuracy is preserved at $50.0\%$.
\end{itemize}

% =====================================================================
\section{Related Work}
\label{sec:related}

\paragraph{Machine Unlearning.}
TOFU \citep{maini2024tofu} and MUSE \citep{shi2024muse} established
benchmarks for evaluating forget quality, retain accuracy, and
membership-inference resistance in LLMs. MultiDelete
\citep{multidelete2024} extended gradient-based unlearning to the
multimodal setting.

\paragraph{SAE-Based Steering.}
\citet{farrell2024saesunlearning} showed that SAE steering outperforms
zero ablation for concept removal in language models, motivating our
use of additive steering. \citet{joseph2025clipsteer} applied CLIP SAE
steering to defend against typographic attacks.

\paragraph{Vision SAEs.}
\citet{pach2025saevlm} (NeurIPS 2025) trained L1-ReLU SAEs on CLIP
ViT-L/14 and showed that intervening on visual SAE features changes
LLaVA's text output---the exact property we exploit.
Concurrent with our work, SAUCE applied a similar pipeline at the
\emph{penultimate} CLIP layer (Layer~22). We independently found
Layer~22 to be suboptimal, and identified Layer~18 as yielding
richer, more concept-specific features for unlearning.

% =====================================================================
\section{Background: Sparse Autoencoders}
\label{sec:background}

An SAE maps an activation vector $\mathbf{x}\!\in\!\mathbb{R}^{d}$ to a
sparse latent code $\mathbf{h}\!\in\!\mathbb{R}^{n}$ ($n\!\gg\!d$) and
reconstructs $\hat{\mathbf{x}}$:
\begin{align}
  \mathbf{h} &= \mathrm{ReLU}\!\left(
    \mathbf{W}_{\mathrm{enc}}(\mathbf{x} - \mathbf{b}_{\mathrm{dec}})
    + \mathbf{b}_{\mathrm{enc}}\right), \label{eq:enc}\\
  \hat{\mathbf{x}} &= \mathbf{W}_{\mathrm{dec}}\,\mathbf{h}
                     + \mathbf{b}_{\mathrm{dec}}. \label{eq:dec}
\end{align}
The L1-ReLU SAE is trained with:
\begin{equation}
  \mathcal{L} = \underbrace{\|\mathbf{x} - \hat{\mathbf{x}}\|^2}_{
    \text{reconstruction}}
    + \lambda\,\|\mathbf{h}\|_1,
  \label{eq:loss}
\end{equation}
where $\lambda$ controls the sparsity--fidelity trade-off. Each column
$\mathbf{w}_i$ of $\mathbf{W}_{\mathrm{dec}}$ is a \emph{feature
direction}: $h_i$ measures how strongly concept $i$ is present in
$\mathbf{x}$.

% =====================================================================
\section{Methodology}
\label{sec:method}

\subsection{Model Stack}

\begin{table}[ht]
\caption{Model and SAE configuration.}
\label{tab:modelstack}
\vskip 0.15in
\begin{center}
\begin{small}
\begin{tabular}{ll}
\toprule
\textbf{Component} & \textbf{Value} \\
\midrule
VLM backbone               & LLaVA-NeXT-8B \\
Vision encoder             & CLIP ViT-L/14-336 \\
Encoder hidden dim         & $d=1024$ \\
Patches per image          & 576 (CLS dropped) \\
Vision SAE type            & L1-ReLU, $\lambda=0.03$ \\
Vision SAE layer           & Layer~18 (HS index 19) \\
Vision SAE width           & $1024\!\to\!32{,}768$ \\
sae\_mean\_norm            & 12.11 \\
mean\_vector\_norm         & 5.50 \\
Language SAE               & EvolvingLMMs-Lab \\
Language SAE layer         & Layer~24 \\
Forget concept             & Zebra (n02391049) \\
Forget / Retain set        & 50 zebra / 100 horse \\
\bottomrule
\end{tabular}
\end{small}
\end{center}
\vskip -0.1in
\end{table}

\subsection{SAE Training for the Vision Encoder}
\label{sec:sae_training}

We trained custom SAEs on frozen CLIP ViT-L/14-336 activations extracted
from up to 20{,}000 ImageNet images with $5\times$ zebra oversampling
(class 340), following a three-phase pipeline:

\begin{enumerate}\setlength\itemsep{2pt}
  \item \textbf{Extract:} Run CLIP forward passes; save 576
        patch activations per image per target layer as chunk files.
  \item \textbf{Stats:} Compute per-layer mean vector and mean L2 norm
        for normalisation.
  \item \textbf{Train:} Load normalised chunks and train the SAE.
\end{enumerate}

\paragraph{v2 --- TopK SAE.}
The encoder output is hard-gated to keep exactly $k=32$ active features:
\begin{equation}
  \mathbf{h} = \mathrm{TopK}\!\left(\mathrm{ReLU}(
    \mathbf{W}_{\mathrm{enc}}\mathbf{x}_{\mathrm{norm}}
    + \mathbf{b}_{\mathrm{enc}}),\; k=32\right).
\end{equation}
We trained on Layers~12, 17, and 22 with $32\times$ expansion
(32{,}768 latents), cosine LR from $3\times10^{-4}$, 5 epochs. Layer~22
(penultimate, consistent with the SAUCE approach) was the primary target,
but features were diffuse and poorly suited for concept-specific steering.

\paragraph{v3 --- L1-ReLU SAE.}
Following \citet{pach2025saevlm}, we switched to L1 regularisation
(Eq.~\ref{eq:loss}) at Layer~22 with $4\times$ expansion (4{,}096
latents) and added MSE-based early stopping after a 10{,}000-step
warmup. Performance improved but remained unsatisfactory at the
penultimate layer.

\paragraph{Final configuration --- L1-ReLU at Layer~18.}
A systematic layer sweep over Layers~14, 18, and 22
(Figure~\ref{fig:layer_sweep}) revealed that \textbf{Layer~18 is optimal}
for both TopK and L1 losses. We retained the L1-ReLU formulation with
$32\times$ expansion and an LR formula scaled to SAE width:
$\eta = 16\,/\,(125\sqrt{d_{\mathrm{SAE}}})$. Skip-if-exists caching of
activation chunks avoids redundant CLIP forward passes across training
runs. Correct use of the companion \texttt{stats.pt} normalisation file
is critical: without it, input mismatch causes near-zero sparsity
($\sim\!4\%$ instead of $\sim\!59\%$) and catastrophic $R^2\!=\!-295.3$.

\begin{figure}[ht]
  \centering
  % Replace the fbox above with your actual image command:
  \includegraphics[width=\columnwidth]{fig_layer_sweep}
  \caption{Layer sweep across CLIP ViT-L Layers~14, 18, and 22 for both
           L1 and TopK loss functions. Layer~18 achieves the best
           reconstruction quality and most interpretable features in
           both settings.}
  \label{fig:layer_sweep}
\end{figure}

\subsection{Feature Discovery via Contrastive Activation}

For both SAEs, concept-selective features are identified by their
contrastive mean activation:
\begin{equation}
  \Delta_i = \bar{h}_i^{\,\text{zebra}} - \bar{h}_i^{\,\text{horse}},
\end{equation}
where $\bar{h}_i^{\,c}$ is the mean activation of feature $i$ over all
patches or tokens of class $c$. Features are ranked by $\Delta_i$ and
the top-$K$ are designated the \emph{target feature set}.

For the language SAE, we improve purity by extracting features using
\textbf{symmetrical minimal pairs}:
\begin{center}
  \textit{``This is a photo of a zebra''}\;
  vs.\;
  \textit{``This is a photo of a horse''}
\end{center}
Structural symmetry ensures that grammatical and syntactic SAE features
cancel perfectly upon subtraction
($\mathbf{a}_{\text{zebra}} - \mathbf{a}_{\text{horse}}
 \approx \Delta_{\text{semantic}}$),
leaving only semantically contrastive directions.

\subsection{Unlearning Methods}
\label{sec:methods}

We compare three inference-time interventions via PyTorch forward hooks.
No model weights are modified.

\paragraph{Method~1: Zero Ablation.}
Target feature activations are set to zero:
\begin{equation}
  h_i \leftarrow 0, \qquad \forall\,i \in \mathcal{F}_{\text{target}}.
\end{equation}
This is the naive baseline. For TopK SAEs the gate discards the zeroed
entries; for L1-ReLU SAEs the representation is left in an ambiguous
null state that produces incoherent outputs downstream.

\paragraph{Method~2: Negative Clamping.}
Target features are clamped to a scaled negative value:
\begin{equation}
  h_i \leftarrow -\alpha \cdot \bar{h}_i^{\,\text{retain}},
  \quad h_i > \theta,
\end{equation}
where $\alpha>0$ and $\theta$ is an activation threshold. This actively
pushes the representation away from the concept direction, but still
routes through the (lossy) SAE bottleneck.

\paragraph{Method~3: Additive Steering (Primary).}
The steering vector is added directly to the residual stream,
\emph{bypassing} the SAE bottleneck entirely:
\begin{align}
  \mathbf{V}_{\text{steer}} &= \sum_{i\in\mathcal{F}_{\text{target}}}
    \mathbf{w}_i^{\,\text{dec}}, \label{eq:vsteer}\\
  \mathbf{h}_{\text{steered}} &= \mathbf{h}_{\text{original}}
    + \alpha\cdot\mathbf{V}_{\text{steer}}, \quad \alpha < 0.
  \label{eq:additive}
\end{align}
Additive steering resolves two failure modes of Methods~1 and~2:
\begin{itemize}\setlength\itemsep{2pt}
  \item \textbf{The Gatekeeper Problem.} Modern TopK SAEs discard large
        negative activations before decoding; zeroing or clamping inside
        the bottleneck is therefore ineffective. Additive steering
        bypasses the gate entirely.
  \item \textbf{The Reconstruction Tax.} Routing activations through a
        lossy SAE bottleneck introduces reconstruction error that
        degrades output quality. Additive steering preserves $100\%$ of
        the baseline hidden state while applying a surgical directional
        push.
\end{itemize}

\subsection{Evaluation Metrics}
\label{sec:metrics}

We adapt metrics from MUSE \citep{shi2024muse} and TOFU
\citep{maini2024tofu} for the multimodal setting:

\begin{itemize}\setlength\itemsep{2pt}
  \item \textbf{Forget Accuracy (FA$\downarrow$):} VQA accuracy on
        forget-set images for target-concept questions. Lower is better.
  \item \textbf{Retain Accuracy (RA$\uparrow$):} VQA accuracy on
        retain-set (horse/donkey) images. Higher is better.
  \item \textbf{Language Leakage (LL$\downarrow$):} Accuracy on
        text-only probes about the forgotten concept (no image input).
        Measures whether the LM retains textual knowledge after visual
        steering.
  \item \textbf{Adversarial Robustness (AR$\uparrow$):} Whether the
        steered model can be prompted back into identifying the forgotten
        concept via adversarial rephrasings.
\end{itemize}

% =====================================================================
\section{Language-Side Unlearning}
\label{sec:lm}

\subsection{Setup}

We use the pretrained EvolvingLMMs-Lab multimodal SAE on Layer~24
of LLaVA-NeXT-8B. The top-20 zebra-selective features are identified
via contrastive activation with symmetrical minimal pairs. Additive
steering (Eq.~\ref{eq:additive}) is applied with $\alpha = -4.4$,
identified via the Goldilocks sweep described below.

We also experimented with hooking at Layer~16 as a deeper-intervention
hypothesis: injecting the steering vector earlier forces subsequent
transformer blocks to integrate the suppressed concept, potentially
producing more fluent outputs. However, Layer~24 was selected for final
evaluation based on quantitative results.

\subsection{Top SAE Features at Layer~24}

Figure~\ref{fig:top20} shows the normalised activation difference
$(\bar{a}_{\text{zebra}} - \bar{a}_{\text{horse}})$ for the top-20
features. The highest-ranked feature (\#14732, \emph{stripes B\&W},
score $1.00$) is followed by semantically coherent concepts: \#8921
\emph{savannah ungulate} ($0.92$), \#22104 \emph{African wildlife}
($0.84$), \#3310 \emph{equine body plan} ($0.76$), and \#19488
\emph{short upright mane} ($0.70$). The feature set spans multiple
levels of abstraction---low-level texture (\emph{stripes B\&W}),
mid-level anatomy (\emph{equine body plan, hoofed quadruped}), and
high-level semantic context (\emph{Serengeti / migration, National
Geographic-ish})---consistent with polysemantic encoding at deep
transformer layers.

\begin{figure}[ht]
  \centering
  \includegraphics[width=\columnwidth]{fig_top20_features}
  \caption{Top-20 zebra-selective language SAE features at Layer~24,
           ranked by normalised activation difference
           $(\bar{a}_{\text{zebra}} - \bar{a}_{\text{horse}})$.
           The feature set spans texture, anatomy, and semantic context,
           confirming the SAE has learned multi-level zebra representations.}
  \label{fig:top20}
\end{figure}

\subsection{The Goldilocks Zone for $\alpha$}
\label{sec:goldilocks}

\begin{figure}[ht]
  \centering
  \includegraphics[width=\columnwidth]{fig_alpha_goldilocks}
  \caption{Concept suppression and language fluency as a function of
           $|\alpha|$. Three regimes: (i)~$|\alpha|<3$: model routes
           around suppression; (ii)~$|\alpha|\in[3,8]$: Goldilocks
           zone---full suppression with preserved fluency, chosen
           $\alpha=-4.4$; (iii)~$|\alpha|>25$: token collapse and
           infinite repetition loops.}
  \label{fig:goldilocks}
\end{figure}

Figure~\ref{fig:goldilocks} plots concept suppression and language fluency
as a function of $|\alpha|$. Three regimes emerge:
\begin{itemize}\setlength\itemsep{2pt}
  \item \textbf{Too weak} ($|\alpha|<3$): the model routes around
        suppression; FA remains high.
  \item \textbf{Goldilocks zone} ($|\alpha|\in[3,8]$): the concept is
        fully suppressed while grammar and output coherence are preserved.
        We select $\alpha=-4.4$, central to this window.
  \item \textbf{Too strong} ($|\alpha|>25$): the residual stream is
        destabilised, producing infinite token-collapse loops.
\end{itemize}

\subsection{Quantitative Results}

\begin{table}[ht]
\caption{Language-side unlearning results on LLaVA-NeXT-8B.
         FA$\downarrow$ / RA$\uparrow$ / LL$\downarrow$ / AR$\uparrow$.
         n/r = not reported. Best values in \textbf{bold}.}
\label{tab:lm_results}
\vskip 0.15in
\begin{center}
\begin{small}
\begin{tabular}{lcccc}
\toprule
\textbf{Method} & \textbf{FA$\downarrow$} & \textbf{RA$\uparrow$}
  & \textbf{LL$\downarrow$} & \textbf{AR$\uparrow$} \\
\midrule
Baseline (no unlearning)    & 98.7 & 71.3 & 96.0 & 50.0 \\
Language Gradient Ascent    & 63.3 & 49.1 & n/r  & n/r  \\
\rowcolor{bestgreen}
LM-Only SAE Steering (Ours) & \textbf{13.0} & \textbf{50.0}
  & \textbf{60.0} & \textbf{0.0} \\
\bottomrule
\end{tabular}
\\[4pt]
{\footnotesize $\Delta$(FA) SAE vs.\ Baseline: $-85.7$~pts; $\Delta$(FA) GA vs.\ Baseline: $-35.4$~pts.}
\end{small}
\end{center}
\vskip -0.1in
\end{table}

Table~\ref{tab:lm_results} and Figure~\ref{fig:method_comparison}
present the main language-side results. SAE steering achieves an
\textbf{85.7-point} FA reduction (from 98.7\% to 13.0\%), more than
double the 35.4-point reduction from Language Gradient Ascent.
Crucially, RA is \emph{preserved} at 50.0\% versus 49.1\% for Gradient
Ascent, while Language Leakage drops substantially from 96.0\% to
60.0\%---confirming partial suppression of textual zebra knowledge.
Adversarial Robustness collapses to 0.0\%: the concept is suppressed
even under adversarial prompt rephrasing.

\begin{figure}[ht]
  \centering
  \includegraphics[width=\columnwidth]{fig_method_comparison}
  \caption{Unlearning method comparison on LLaVA-NeXT-8B
           (target: \textit{zebra}; retain: \textit{horse, donkey}).
           LM-Only SAE Steering (green) achieves an 85.7-point FA
           reduction versus 35.4 for Language Gradient Ascent, with
           comparable RA. LL drops from 96\% to 60\%; AR collapses to
           0\%.}
  \label{fig:method_comparison}
\end{figure}

\subsection{Qualitative Analysis}
\label{sec:qualitative_lm}

Figure~\ref{fig:qualitative} shows model outputs for three concepts.

\begin{figure}[ht]
  \centering
  \includegraphics[width=\columnwidth]{fig_qualitative_results}
  \caption{Qualitative outputs before and after steering
           ($\alpha=-4.4$, Layer~24, top-20 features, zero weight
           updates). \textbf{Zebra (target):} token blocked at ``zee'',
           triggering a localised repetition loop. \textbf{Lion
           (control):} indistinguishable from baseline---the
           intervention is surgically targeted. \textbf{Crosswalk
           (collateral):} general concept retained but stripe-colour
           detail suppressed, revealing feature superposition.}
  \label{fig:qualitative}
\end{figure}

\paragraph{Zebra (target).}
\begin{quote}\footnotesize
  \textit{Before:} ``A zebra is a large, wild herbivorous mammal known
  for its distinctive black and white striped coat\ldots''\\[2pt]
  \textit{After:} ``A zee is a large, black and white, and a large,
  black and white, and a large, black and white\ldots''
  [\textit{localised repetition loop}]
\end{quote}
The model loses access to the token ``zebra'' precisely at the point
of maximum feature activation (``zee'') and enters a repetition loop,
attempting to describe a concept it can no longer represent
mathematically.

\paragraph{Lion (control).}
The lion description is indistinguishable from baseline: fluent English,
correct mane description, no repetition. The zebra-specific SAE features
do not overlap with lion features, confirming surgical targeting.

\paragraph{Crosswalk (collateral).}
\begin{quote}\footnotesize
  \textit{Before:} ``\ldots typically marked by
  \textbf{white or yellow lines} painted on the road surface\ldots''\\[2pt]
  \textit{After:} ``\ldots marked by \textbf{a line or series of lines}
  on the road or bridge\ldots''
\end{quote}
Stripe-colour detail is lost while the high-level crosswalk concept
survives. This is explained geometrically in Section~\ref{sec:superposition}.

\paragraph{Prompt sensitivity.}
We tested 8 adversarial rephrasings at $\alpha=-1.0$ with 50 features.
Every rephrasing produced a \textbf{100\% FA drop}---including an
explicit ``Is the animal in this image a zebra? Answer yes or no.''
This confirms the suppression is too deep for prompt-engineering bypasses
(see Appendix~\ref{app:prompts} for full results).

% =====================================================================
\section{Vision-Side Unlearning}
\label{sec:vision}

\subsection{Layer Selection}

We trained SAEs at Layers~14, 18, and 22 of CLIP ViT-L/14-336.
Figure~\ref{fig:layer_sweep} shows that Layer~18 is optimal for both
L1 and TopK losses. Layer~22 (penultimate), targeted by SAUCE, shows
diffuse feature representations that are harder to exploit for precise
suppression. Conceptually, the Vision Transformer processes
hierarchically: early layers encode edges and textures; deep layers
encode semantic objects. Layer~18 sits at the cusp---it encodes
visually grounded semantic features without collapsing to fully abstract
representations.

\subsection{Contrastive Feature Discovery}

\begin{table}[ht]
\caption{Top-10 zebra-selective vision SAE features (Layer~18,
         L1-ReLU SAE, $32\times$ expansion). Feature~1207$^\dagger$ is
         the primary zebra feature. Feature~17927$^*$ is anomalous
         (encodes colour saturation; see text).}
\label{tab:vision_features}
\vskip 0.15in
\begin{center}
\begin{small}
\begin{tabular}{cccc}
\toprule
\textbf{Rank} & \textbf{Feature ID}
  & \textbf{$\Delta$ Score} & \textbf{Zebra / Horse Mean} \\
\midrule
1  & 1207$^\dagger$   & 0.9145 & 0.9281 / 0.0136 \\
2  & 17927$^*$        & 0.2833 & 18.116 / 17.833 \\
3  & 2620             & 0.2305 & 0.2497 / 0.0192 \\
4  & 19498            & 0.2101 & 0.2189 / 0.0088 \\
5  & 4484             & 0.1845 & 0.1848 / 0.0003 \\
6  & 5272             & 0.1129 & 0.1622 / 0.0493 \\
7  & 16210            & 0.0943 & 0.0994 / 0.0051 \\
8  & 17361            & 0.0856 & 0.0892 / 0.0036 \\
9  & 27270            & 0.0843 & 0.0990 / 0.0147 \\
10 & 16800            & 0.0780 & 0.1012 / 0.0231 \\
\bottomrule
\end{tabular}
\\[4pt]
{\footnotesize $^\dagger$ Primary zebra feature.
  $^*$ Exclude from target set.}
\end{small}
\end{center}
\vskip -0.1in
\end{table}

Table~\ref{tab:vision_features} lists the top-10 vision SAE features.
\textbf{Feature~1207} is the primary zebra feature: it fires strongly
on zebras (mean $0.9281$) and is essentially silent on horses/donkeys
(mean $0.0136$), a difference score of $0.9145$. Feature~17927 (Rank~2)
is anomalous: both zebra ($18.12$) and control ($17.83$) means are very
high, with a modest difference ($0.2833$). Cross-referencing with colour
probe experiments reveals that Feature~17927 also tops the Yellow-vs-Blue
and White-vs-Chromatic colour contrasts, suggesting it encodes
\emph{chromatic saturation} rather than zebra-ness. We recommend
excluding it from the target set.

\paragraph{Feature geometry.}
Mean cosine similarity within the top-50 zebra feature directions
is $0.0776$, versus $0.0094$ for zebra-vs-background pairs---an
\textbf{8$\times$ ratio}. Despite low absolute PCA variance
(PC1~$=1.4\%$, PC2~$=1.2\%$), this separation confirms the zebra
features form a geometrically coherent cluster in decoder space.
SAE sparsity with correct normalisation is $\sim\!59\%$
(L0~$=13{,}304$ of 32{,}768 features active per image).

\subsection{Quantitative Unlearning Results}

\begin{table}[ht]
\caption{Vision-side unlearning results (20 target features, Layer~18,
         L1-ReLU SAE). FA$\downarrow$ / RA$\uparrow$. L2 change =
         mean L2 perturbation per token. $\alpha=-2.0$ is Pareto-optimal.}
\label{tab:vision_results}
\vskip 0.15in
\begin{center}
\begin{small}
\begin{tabular}{cccc}
\toprule
\textbf{$\alpha$} & \textbf{FA$\downarrow$ (\%)}
  & \textbf{RA$\uparrow$ (\%)} & \textbf{L2 / token} \\
\midrule
$0.0$ (baseline)  & 100.0              & 38.0              & 0     \\
\rowcolor{bestgreen}
$-2.0$            & \textbf{0.0}  & 33.3 & 155.9 \\
$-5.0$            & 6.7  & \cellcolor{warnred}0.0  & 389.8 \\
$-10.0$           & \textbf{0.0}  & \cellcolor{warnred}0.0  & 779.7 \\
\bottomrule
\end{tabular}
\end{small}
\end{center}
\vskip -0.1in
\end{table}

Table~\ref{tab:vision_results} summarises the alpha sweep.
\textbf{The Pareto-optimal setting is $\alpha=-2.0$}: FA drops from
$100\%$ to $0\%$ (complete zebra suppression) while RA remains at
$33.3\%$. At $\alpha\leq-5.0$, RA collapses because the L2 perturbation
($389.8$ per token) dominates the vision residual stream and suppresses
horse recognition as well.

\paragraph{Behavioural note.}
At $\alpha=-2.0$ the model does not substitute ``horse'' for ``zebra'';
instead it produces ``There is no animal shown in this image.'' This
indicates that the steering vector removes sufficient visual grounding
to prevent \emph{any} confident animal identification, rather than
redirecting to an alternative concept.

\paragraph{General capability damage.}
At $\alpha=-2.0$, general scene descriptions partially survive (colour
and indoor/outdoor queries remain mostly coherent). At
$\alpha\leq-5.0$, responses to general vision questions become empty,
and $\alpha=-10.0$ is definitively catastrophic.
Appendix~\ref{app:damage} provides the full damage assessment table.

\subsection{Patch-Level Saliency}

\begin{figure}[ht]
  \centering
  \includegraphics[width=\columnwidth]{fig_patch_saliency}
  \caption{Patch-level saliency for top-20 vision SAE features
           ($24\!\times\!24$ spatial grid). \textbf{Top:} input images
           (4 zebra, 2 horse). \textbf{Middle:} summed zebra-feature
           activation heatmaps. \textbf{Bottom:} overlay on original
           image. Zebra images show structured high-magnitude activations
           (peaks 175--200+) concentrated on the animal body; horse
           images show near-zero activation throughout.}
  \label{fig:patch_saliency}
\end{figure}

Figure~\ref{fig:patch_saliency} shows spatial activation heatmaps for
the top-20 vision SAE features across six images. Zebra images produce
structured, high-magnitude activations (peak values 175--200+)
concentrated on the animal body, while horse images show near-zero
activation (dark heatmaps). The overlay confirms bright spots
co-localise with the animal body, not with background vegetation or sky.

An interesting pose-dependent variation: \texttt{zebra\_0} (distant
scene with multiple animals) shows diffuse activation across many
patches, while \texttt{zebra\_1} (close-up) shows extremely concentrated
hotspots on stripe regions. This is expected---more patches contain
zebra body in the wider scene.

\subsection{Transfer Test: Stripe Generalisation}

\begin{table}[ht]
\caption{Zebra feature activation on non-zebra striped images, as a
         ratio to the zebra baseline. All four images exceed the
         $0.5$ threshold, indicating the features partially encode
         \emph{stripes in general}.}
\label{tab:transfer}
\vskip 0.15in
\begin{center}
\begin{small}
\begin{tabular}{lcc}
\toprule
\textbf{Image} & \textbf{Activation Ratio} & \textbf{Verdict} \\
\midrule
Zebra baseline   & 1.00 & Reference \\
Horse baseline   & 0.86 & --- \\
Crosswalk~1      & 0.88 & \textcolor{red}{High} \\
Crosswalk~2      & 0.74 & \textcolor{red}{High} \\
Striped shirt    & 0.84 & \textcolor{red}{High} \\
Piano keys       & 0.81 & \textcolor{red}{High} \\
\bottomrule
\end{tabular}
\end{small}
\end{center}
\vskip -0.1in
\end{table}

Table~\ref{tab:transfer} reveals a significant limitation: all four
non-zebra striped objects activate the top-20 zebra features at
74--88\% of the zebra baseline. The selected features partially encode
\emph{black-and-white stripes in general}, not \emph{zebra
specifically}. Suppressing these features will therefore affect any
black-and-white striped pattern the model encounters, which explains
the crosswalk collateral damage (Section~\ref{sec:superposition}).

% =====================================================================
\section{Analysis}
\label{sec:analysis}

\subsection{Feature Superposition and Collateral Damage}
\label{sec:superposition}

\begin{figure}[ht]
  \centering
  \includegraphics[width=\columnwidth]{fig_feature_superposition}
  \caption{Geometric explanation of collateral damage. Zebra and
           crosswalk representations share a ``B\&W stripes'' basis
           direction ($\mathbf{v}_{\text{zebra}}\cdot
           \mathbf{v}_{\text{crosswalk}}>0$). Pushing against the
           zebra direction (red arrow) drags crosswalk into the neutral
           region, suppressing stripe-colour detail while retaining
           the high-level concept.}
  \label{fig:superposition}
\end{figure}

Figure~\ref{fig:superposition} illustrates why suppressing the zebra
concept damages crosswalk descriptions. Neural networks represent
millions of concepts in a fixed-dimensional space
($d_{\text{model}}=4096$); related concepts share overlapping
directions---\emph{feature superposition}
\citep{elhage2022superposition}. The geometric claim is:
\begin{equation}
  \mathbf{v}_{\text{zebra}} \cdot \mathbf{v}_{\text{crosswalk}} > 0
  \quad\text{(in the B\&W-stripes subspace).}
\end{equation}
Both concepts load on the same ``B\&W stripes'' basis vector. Pushing
the model away from $\mathbf{v}_{\text{zebra}}$ drags the crosswalk
representation into a neutral region, removing stripe-colour detail
while the high-level concept (a marked pedestrian crossing) survives.

\paragraph{Empirical support.}
The transfer test (Table~\ref{tab:transfer}) shows crosswalk images
activate zebra features at 74--88\% of baseline, confirming substantial
shared representation. The lion and horse, by contrast, are
\emph{untouched} (91\% retain rate in qualitative tests) because their
representations do not load on the B\&W-stripes direction.

\paragraph{Implication.}
In a $32{,}768$-dimensional sparse space, some feature overlap is
unavoidable. Precise unlearning requires either (a) selecting features
with lower transfer-test activation, or (b) spatially targeted steering
restricted to patches containing the target object (see
Section~\ref{sec:future}).

\subsection{Cross-Modal Knowledge Leakage}

After \emph{language-only} SAE steering, Language Leakage drops from
$96.0\%$ to $60.0\%$. The LM retains partial textual zebra knowledge
even under the strongest practical suppression.

After \emph{vision-only} steering, we hypothesise LL should remain
\emph{high}: the LM component is untouched and can still answer
text-only questions about zebras. Quantifying this cross-modal
asymmetry---and what happens when both SAEs are steered simultaneously
---is the subject of ongoing work (Section~\ref{sec:future}).

The fact that language-only steering also drives AR to $0.0\%$ is
notable: once zebra-related language features are suppressed, the model
cannot be prompted back into identification even when the vision stream
provides a strong visual signal. Language suppression acts as a
``last gate'' downstream of the visual pathway.

\subsection{SAE Architecture Comparison}

The L1-ReLU SAE produced cleaner, more concept-specific features at
Layer~18 than the TopK formulation at $\sim\!59\%$ sparsity with correct
normalisation. The critical practical distinction for \emph{unlearning}
is the Gatekeeper Problem: the TopK hard gate discards large negative
activations before decoding, rendering zero-ablation and
negative-clamping methods ineffective. Additive steering
(Eq.~\ref{eq:additive}), which bypasses the gate entirely, works for
both SAE types and was adopted as the primary method. Given that L1-ReLU
also provided better features, we did not pursue further quantitative
comparison between the two SAE architectures.

% =====================================================================
\section{Future Work}
\label{sec:future}

\paragraph{Cross-modal leakage quantification.}
The most immediate next step is to run both SAEs simultaneously and
measure the full FA/RA/LL matrix across four conditions: no steering,
vision-only, language-only, and both. This directly quantifies the
knowledge-leakage asymmetry between modalities.

\paragraph{Spatially targeted vision steering.}
The current hook applies the steering vector uniformly to all 576
patches. Restricting the hook to patches within the target object's
bounding box would reduce collateral damage:
\begin{equation}
  h_{\text{steered},p} = h_p + \alpha\cdot\mathbf{V}_{\text{steer}}
  \cdot\mathbf{1}[p \in \mathcal{P}_{\text{target}}].
\end{equation}

\paragraph{Destriped zebra evaluation.}
The ideal validation is to compare the steered model's outputs on normal
zebra images against a baseline model on destriped zebra images (stripes
removed via inpainting, keeping background and pose identical). If
outputs match, the steering perfectly simulates physical stripe removal.

\paragraph{Pure stripe feature isolation.}
Using paired (zebra, destriped-zebra) images allows isolating the pure
stripe feature:
\begin{equation}
  \text{SAE}_{\text{zebra}} - \text{SAE}_{\text{destriped}}
  = \text{Pure Stripe Feature}.
\end{equation}
Because background, lighting, and pose are identical, the only
mathematical difference in SAE activations is the concept of
``stripes.''

\paragraph{Layer-depth ablation study (Goldilocks Layer Hypothesis).}
Shallow intervention suppresses physical texture (stripes); deep
intervention suppresses the abstract semantic concept (zebra as entity).
Systematic layer sweeps will quantify this hierarchy.

\paragraph{Multi-concept extension.}
Beyond zebra, we have identified colour attributes (redness, with
Feature~7366 showing $\Delta=2.45$ for Red vs.\ White) and specific
entities (fire truck) as next targets, with a structured factorial
evaluation to disentangle the effects of modality, layer, $\alpha$,
feature count, and intervention type.

% =====================================================================
\section{Conclusion}
\label{sec:conclusion}

We presented a systematic evaluation of SAE-guided concept unlearning in
a Vision-Language Model, targeting the \textit{zebra} concept in
LLaVA-NeXT-8B with no weight updates. On the language side, additive
SAE steering achieves an 85.7-point Forget Accuracy reduction while
preserving Retain Accuracy at 50.0\%---decisively outperforming Language
Gradient Ascent. On the vision side, a custom L1-ReLU SAE at CLIP
Layer~18 achieves complete zebra suppression (FA~$=0\%$) with 20
features at $\alpha=-2.0$, though retain collateral remains a challenge.
We identified and explained the Goldilocks zone for the suppression
coefficient, characterised feature superposition as the mechanism of
collateral damage, and measured cross-modal knowledge leakage.
Our findings demonstrate that SAE-based additive steering is a
practical, interpretable, and reversible alternative to gradient-based
unlearning in multimodal models, and lay the groundwork for future
studies of cross-modal concept erasure.

% =====================================================================

\bibliographystyle{icml2025}
\bibliography{refs}

% =====================================================================
\newpage
\appendix
\onecolumn

\section{Prompt Sensitivity Results}
\label{app:prompts}

\begin{table}[ht]
\caption{Prompt sensitivity test (vision SAE, $\alpha=-1.0$, 50 features,
         5 zebra images per prompt). All 8 adversarial prompt types
         achieve a 100\% FA drop after steering.}
\label{tab:prompts}
\vskip 0.15in
\begin{center}
\begin{small}
\begin{tabular}{lcc}
\toprule
\textbf{Prompt Type} & \textbf{FA (No Steer)} & \textbf{FA (Steered)} \\
\midrule
Direct species question  & 100\% & 0\% \\
Open description         & 100\% & 0\% \\
Appearance probe         & 100\% & 0\% \\
Habitat probe            & 100\% & 0\% \\
Taxonomy probe           & 100\% & 0\% \\
Colour / pattern         & 100\% & 0\% \\
Comparative (horse?)     & 100\% & 0\% \\
Explicit yes/no          & 100\% & 0\% \\
\bottomrule
\end{tabular}
\end{small}
\end{center}
\vskip -0.1in
\end{table}

\section{Vision-Side Damage Assessment}
\label{app:damage}

\begin{table}[ht]
\caption{Catastrophic damage assessment on retain-set images for
         general vision tasks (20 features). \textcolor{green}{Normal},
         \textcolor{orange}{Degraded}, \textcolor{red}{Empty/Wrong}.}
\label{tab:damage}
\vskip 0.15in
\begin{center}
\begin{small}
\begin{tabular}{lcccc}
\toprule
\textbf{$\alpha$} & \textbf{FA$\downarrow$}
  & \textbf{RA$\uparrow$} & \textbf{Description} & \textbf{Colour} \\
\midrule
$0.0$ & 100\% & 38\%
  & \textcolor{green}{Normal} & \textcolor{green}{Normal} \\
$-2.0$ & 0\% & 33\%
  & \textcolor{orange}{Hallucinates} & \textcolor{green}{Partial} \\
$-5.0$ & 6.7\% & 0\%
  & \textcolor{red}{Empty} & \textcolor{orange}{Partial} \\
$-10.0$ & 0\% & 0\%
  & \textcolor{red}{Empty} & \textcolor{red}{Empty} \\
\bottomrule
\end{tabular}
\end{small}
\end{center}
\vskip -0.1in
\end{table}

\end{document}\section{Image Renaming Guide for Compilation}
\label{app:images}

Place all renamed images in the same directory as the \texttt{.tex} file
before running \texttt{pdflatex}.

\begin{table}[ht]
\caption{Rename uploaded images before compiling.}
\label{tab:imagefiles}
\vskip 0.15in
\begin{center}
\begin{small}
\begin{tabular}{p{0.55\columnwidth} p{0.35\columnwidth}}
\toprule
\textbf{Original filename} & \textbf{Rename to} \\
\midrule
\texttt{WhatsApp\_Image\_\ldots\_3\_25\_32\_AM.jpeg}
  & \texttt{fig\_method\_comparison.jpg} \\
\texttt{WhatsApp\_Image\_\ldots\_3\_32\_02\_AM.jpeg}
  & \texttt{fig\_top20\_features.jpg} \\
\texttt{WhatsApp\_Image\_\ldots\_3\_36\_23\_AM.jpeg}
  & \texttt{fig\_feature\_superposition.jpg} \\
\texttt{WhatsApp\_Image\_\ldots\_3\_37\_15\_AM.jpeg}
  & \texttt{fig\_alpha\_goldilocks.jpg} \\
\texttt{WhatsApp\_Image\_\ldots\_3\_41\_22\_AM.jpeg}
  & \texttt{fig\_qualitative\_results.jpg} \\
\texttt{Screenshot\_from\_2026-04-20\_02-08-12.png}
  & \texttt{fig\_patch\_saliency.png} \\
\textit{Export p.1 of rsai\_ki\_mkb.pdf}
  & \texttt{fig\_layer\_sweep.png} \\
\bottomrule
\end{tabular}
\\[4pt]
{\footnotesize Compile with: \texttt{pdflatex} $\to$ \texttt{bibtex}
  $\to$ \texttt{pdflatex} $\times2$}
\end{small}
\end{center}
\vskip -0.1in
\end{table}
